\section{Определение линейного функционала. Определение сопряжённого пространства}

\begin{definition}[Линейный функционал]
Пусть $V$ --- линейное пространство над полем $F$, $E$ --- базис в $V$.

Отображение $f: V \to F$ называется \textbf{линейным функционалом}, если:
\begin{enumerate}
    \item $\forall x_1, x_2 \in V: f(x_1 + x_2) = f(x_1) + f(x_2)$;
    \item $\forall x \in V,\ \forall \lambda \in F: f(\lambda x) = \lambda f(x)$.
\end{enumerate}
\end{definition}

\begin{definition}[Сопряжённое пространство]
\textbf{Сопряжённое пространство} $V^*$ --- это множество всех линейных функционалов на $V$ с операциями:
\begin{enumerate}
    \item $f_1 = f_2 \iff \forall x \in V: f_1(x) = f_2(x)$;
    \item $f = 0$, если $\forall x \in V: f(x) = 0$;
    \item $(f_1 + f_2)(x) = f_1(x) + f_2(x),\quad \forall x \in V$;
    \item $(\lambda f)(x) = \lambda f(x),\quad \forall x \in V$.
\end{enumerate}
\end{definition}
